\chapter{Discusión}
\label{ch:discusion}

% Integrar los resultados en el marco teórico presentado en la
% Introducción. Comparar con trabajos previos. Limitaciones del trabajo.

En la presente tesis se diseñaron, extrajeron y evaluaron features a nivel de fijación a partir de los mapas generados por el modelo nnELM, y se estudió su relación con la señal EEG registrada durante una tarea de búsqueda visual híbrida. Los resultados muestran que un subconjunto de las variables latentes del modelo presenta correlatos estadísticamente significativos en la actividad cerebral. En escala original, las features con clusters significativos según el criterio TFCE fueron \feat{posterior}, \feat{saliencia\_media}, \feat{saliencia\_pixel} y \feat{similitud\_pixel}. Al incorporar transformaciones logarítmicas, este conjunto se amplió: \feat{similitud\_media}, \feat{ganancia\_informacion} y \feat{gap\_entropia} revelaron efectos que no eran detectables en su escala original, lo que indica que la relación entre estas variables y la señal EEG es esencialmente no lineal. En conjunto, se encontraron correlatos significativos en features pertenecientes a cuatro de los cinco grupos conceptuales definidos: \textbf{información bottom-up}, \textbf{similitud con el target}, \textbf{estado interno del proceso bayesiano} y \textbf{calidad de la decisión}.

Se pueden realizar distintas interpretaciones para las features que no obtuvieron clusters significativos en ninguna de sus versiones evaluadas. Una de las posibilidades es que estas variables capturen procesos cognitivos cuya expresión neural no es lineal ni capturada adecuadamente por el modelo empleado. En ese caso, métodos no lineales de análisis podrían revelar efectos que no se detectaron en esta oportunidad. 

\pendiente{Análisis de otras posibilidades. Por ejemplo, temporalidad.}

Una limitación central de este trabajo es que las features utilizadas como predictores no fueron extraídas directamente del comportamiento del participante sino de las variables latentes de un modelo computacional. El nnELM es una aproximación al proceso de búsqueda humana: si bien replica métricas de comportamiento con alto rendimiento \citep{ruarte2025}, sus mapas internos son construcciones del modelo que pueden diferir del proceso neural real. En consecuencia, la ausencia de correlatos para una feature determinada no implica necesariamente que el proceso cognitivo asociado no tenga representación neural, sino que el modelo y el cerebro pueden estar resolviendo el mismo problema por caminos distintos. Esta limitación subraya la importancia de interpretar los resultados como evidencia de la correspondencia entre el modelo y la actividad cerebral, y no como evidencia directa sobre los mecanismos neurales de la búsqueda.

Un hallazgo relevante que abre una línea de trabajo futura es la presencia de activación significativa en el período previo al onset de la fijación en las features de saliencia. Esta activación  sugiere que la información de saliencia bottom-up podría estar siendo procesada en relación con la preparación del movimiento ocular más que con el procesamiento del estímulo fijo. En trabajos futuros podría abordarse este fenómeno modelando estas variables en relación con el onset de la sacada en lugar del onset de la fijación, siguiendo la lógica del predictor de intercept de sacada ya incluido en el modelo base.

Por otro lado, el presente trabajo evaluó la contribución de cada feature de forma individual, incorporándola al modelo base de a una. Una extensión natural es explorar modelos que combinen múltiples features simultáneamente, lo que permitiría estudiar no solo el aporte incremental de cada variable sino también las interacciones entre ellas.